%% ============================================================
%% Chapitre : Configuration de la Plateforme ST ChatGPT
%% Fichier séparé pour permettre le repositionnement dans le rapport
%% ============================================================

\chapter{Configuration de la Plateforme ST ChatGPT}
\label{chap:configuration}

Ce chapitre détaille la configuration de la persona \textbf{ST GitHub Analyzer} sur la plateforme ST ChatGPT, les justifications de chaque paramètre, et l'évolution incrémentale du Quality Index qui a guidé nos choix.

\section{Architecture de la Plateforme}

La plateforme ST ChatGPT est un service Azure OpenAI interne déployé par STMicroelectronics. Elle permet de créer des \textbf{personas} --- des agents conversationnels spécialisés --- équipés de deux types d'outils :

\begin{itemize}
  \item \textbf{KB (Knowledge Base)} : recherche dans une base de connaissances structurée
  \item \textbf{Persona-as-a-Tool} : appel à un autre agent comme fallback
\end{itemize}

\section{Configuration du Tool KB}

\subsection{Paramètres Retenus}

\begin{longtable}{p{5cm}p{4cm}p{4cm}}
\toprule
\textbf{Paramètre} & \textbf{Valeur choisie} & \textbf{Alternatives testées} \\
\midrule
KB Type & Unstructured & Structured (rejeté) \\
Search Type & \textbf{Hybrid} & Semantic seul, Full-text seul \\
Semantic Threshold & \textbf{0.55} & 0.3, 0.4, 0.7, 0.8 \\
Semantic Max Doc Count & \textbf{20} & 5, 10, 30 \\
Full-text Max Doc Count & \textbf{12} & 5, 8, 20 \\
Modèle LLM & GPT-5.1 & GPT-4o (précédent) \\
\bottomrule
\end{longtable}

\subsection{Justification : Recherche Hybride}

La recherche \textbf{hybride} (Semantic + Full-text) est supérieure aux deux modes isolés dans notre contexte :

\begin{longtable}{p{3cm}p{5cm}p{5cm}}
\toprule
\textbf{Mode} & \textbf{Force} & \textbf{Faiblesse pour STM32} \\
\midrule
Semantic seul & Comprend le sens (``mon SPI marche pas'' → issues SPI) & Rate les noms exacts d'API (\texttt{HAL\_SPI\_Transmit\_DMA}) \\
Full-text seul & Match exact sur les identifiants techniques & Rate les reformulations (``erreur Ethernet'' vs ``LwIP crash'') \\
\textbf{Hybrid} & Combine les deux & Complexité légèrement supérieure \\
\bottomrule
\end{longtable}

\textbf{Résultat mesuré} : La recherche hybride améliore la couverture de +22\% par rapport au semantic seul sur notre jeu de test de 52 questions, car les questions techniques STM32 mélangent du langage naturel et des identifiants de code.

\subsection{Justification : Seuil Sémantique à 0.55}
\label{sec:seuil_semantique}

Le \textbf{seuil sémantique} (Semantic Threshold) est le score minimum de similarité qu'un document doit atteindre pour être retourné par la recherche sémantique. Un score de 1.0 signifie une correspondance parfaite, 0.0 aucune similarité. Ce seuil contrôle directement le compromis entre \textbf{rappel} (ne rien rater) et \textbf{précision} (pas de bruit).

Le seuil de 0.55 a été calibré empiriquement :

\begin{longtable}{cccp{5cm}}
\toprule
\textbf{Seuil} & \textbf{Docs retournés} & \textbf{QI} & \textbf{Analyse} \\
\midrule
0.3 & $\sim$45 & 48 & Trop de bruit, documents non pertinents diluent le contexte \\
0.4 & $\sim$32 & 51 & Meilleur mais encore des faux positifs \\
\textbf{0.55} & $\sim$20 & \textbf{54} & Équilibre optimal signal/bruit \\
0.7 & $\sim$8 & 47 & Trop restrictif, rate des documents pertinents \\
0.8 & $\sim$3 & 39 & Presque aucun document retourné \\
\bottomrule
\end{longtable}

Le seuil 0.55 maximise le Quality Index en retournant suffisamment de documents pour couvrir les questions variées tout en éliminant les documents trop faiblement liés.

\subsection{Justification : 20 Documents Sémantiques + 12 Full-text}
\label{sec:max_documents}

Le paramètre \textbf{Max Doc Count} limite le nombre de documents retournés par chaque mode de recherche. Ce nombre est injecté dans le contexte du LLM, qui doit ensuite synthétiser une réponse à partir de ces documents. Un nombre trop élevé sature la fenêtre de contexte et ralentit la réponse ; un nombre trop bas risque de manquer le document pertinent.

\begin{itemize}
  \item \textbf{20 documents sémantiques} : couvre les reformulations et les questions indirectes. Au-delà de 20, les documents supplémentaires sont rarement pertinents (score $< 0.55$) et ajoutent du bruit au contexte. Ce chiffre a été déterminé en observant que pour 90\% des questions, le document pertinent se trouve dans les 15 premiers résultats sémantiques.
  \item \textbf{12 documents full-text} : assure que les correspondances exactes d'API et de noms de fichiers sont toujours retrouvées. Le chiffre 12 couvre les cas où un même identifiant apparaît dans plusieurs séries MCU (ex: \texttt{HAL\_SPI\_Init} existe dans F4, H5, H7, U5, WL).
\end{itemize}

\textbf{Total de documents par requête} : $20 + 12 = 32$ documents maximum, soit environ 16\,000 tokens de contexte (500 tokens/document en moyenne), ce qui reste confortablement dans la fenêtre de contexte de GPT-5.1 (128K tokens).

\subsection{Définition : Search Type (Hybrid)}
\label{sec:search_type}

Le \textbf{Search Type} détermine l'algorithme de recherche utilisé pour interroger la Knowledge Base. La plateforme ST ChatGPT supporte trois modes :

\begin{enumerate}
  \item \textbf{Semantic} : transforme la requête en vecteur (embedding) et recherche les documents les plus proches dans l'espace vectoriel. Comprend le \textit{sens} de la question mais peut rater les termes exacts.
  \item \textbf{Full-text} : recherche par correspondance exacte de mots (BM25/TF-IDF). Trouve les identifiants exacts (\texttt{HAL\_SPI\_TransmitReceive\_DMA}) mais ne comprend pas les reformulations.
  \item \textbf{Hybrid} : exécute les deux recherches en parallèle et fusionne les résultats (union avec déduplication). Combine la compréhension sémantique et la précision lexicale.
\end{enumerate}

Le mode \textbf{Hybrid} a été retenu car les questions des ingénieurs STM32 mélangent systématiquement du langage naturel (``mon SPI crash au deuxième transfert'') et des identifiants techniques exacts (\texttt{HAL\_SPI\_Transmit\_DMA}, \texttt{NUCLEO-H743ZI}).

\subsection{Tool Description}

La description du tool KB est critique car elle guide le LLM dans sa décision d'appeler ou non le tool :

\begin{quote}
\small\texttt{Retrieve technical knowledge about STM32Cube microcontroller firmware from GitHub repositories. Contains: resolved issues with root cause and fix, HAL/BSP source file documentation, resolver cases (step-by-step troubleshooting), and diagnostic cards (structured fault analysis). Covers series: STM32F4, STM32H5, STM32H7, STM32U5, STM32WL and their driver sub-repos (HAL, CMSIS, BSP). Query this tool when the user asks about: peripheral configuration problems, HAL API behavior, known bugs and workarounds, board-specific issues, firmware package contents, low-power modes, communication protocols, DMA/interrupt issues, or build errors. Use specific technical terms in queries: peripheral names (SPI, I2C, UART, FDCAN, OSPI), HAL function names, error descriptions, board references, and STM32 part numbers.}
\end{quote}

\textbf{Points clés de cette description} :
\begin{enumerate}
  \item \textbf{Énumération explicite du contenu} (issues, files, resolver, diagnostic) → le LLM sait quoi chercher
  \item \textbf{Liste des séries couvertes} → évite les appels pour des séries absentes (L4, G4)
  \item \textbf{Instruction de reformulation} (``Use specific technical terms'') → améliore la qualité des queries envoyées à la KB
  \item \textbf{Exemples de cas d'usage} → le LLM sait \textit{quand} appeler le tool vs répondre directement
\end{enumerate}

\section{Configuration du Tool Alfred (Fallback)}

\subsection{Problème Identifié}

Lors de l'ajout initial d'Alfred comme fallback, le Quality Index a \textbf{chuté de 10 points} (de 54 à 44). L'analyse a montré que la persona appelait Alfred \textbf{systématiquement} au lieu de chercher dans la KB, car la description du tool était trop permissive :

\begin{quote}
\small\texttt{[Description initiale, trop permissive]} ``I can help with STM32 questions. Forward the user's question.''
\end{quote}

\subsection{Fix : Description Restrictive}

La description corrigée force la priorité KB :

\begin{quote}
\small\texttt{I am a fallback assistant for STM32 firmware questions that are NOT covered by the Knowledge Base. The Knowledge Base (KB Unstructured) is ALWAYS the primary source, search it FIRST for every question. Use this tool ONLY when the Knowledge Base returned no relevant result or explicitly lacks coverage on the topic.}
\end{quote}

\subsection{Impact Mesuré}

\begin{longtable}{lcc}
\toprule
\textbf{Configuration} & \textbf{Quality Index} & \textbf{$\Delta$} \\
\midrule
KB seule (sans Alfred) & 54 & baseline \\
KB + Alfred (description permissive) & 44 & $-10$ \\
KB + Alfred (description restrictive) & 52 & $-2$ \\
\bottomrule
\end{longtable}

La leçon : dans une architecture multi-tools, la \textbf{description du tool contrôle le comportement de routage} du LLM. Une description trop ouverte dégrade la performance en court-circuitant la source la plus pertinente.

\section{System Prompt de la Persona}

Le system prompt guide le comportement global de la persona. Les éléments clés :

\begin{enumerate}
  \item \textbf{Identité} : ``You are ST GitHub Analyzer, a technical assistant specialized in STM32Cube MCU firmware.''
  \item \textbf{Priorité KB} : ``ALWAYS search the Knowledge Base FIRST before answering or using any other tool.''
  \item \textbf{Format de réponse} : citations avec liens GitHub, code actionnable, versions spécifiques
  \item \textbf{Stratégie multi-cas} :
  \begin{itemize}
    \item Match direct → citer la source avec URL
    \item Match partiel → diagnostic + workaround candidat
    \item Aucun match → appeler Alfred OU proposer une action concrète
  \end{itemize}
  \item \textbf{Interdiction} : ne jamais dire ``je ne sais pas'' sans avoir cherché dans la KB
\end{enumerate}

\section{Évolution Incrémentale du Quality Index}

Le Quality Index a évolué de manière incrémentale au fil des améliorations du contenu KB et de la configuration. Chaque test mesure la qualité des réponses du chatbot sur un jeu de questions techniques.

\subsection{Chronologie des Tests}

\begin{longtable}{p{2cm}p{5cm}ccp{3.5cm}}
\toprule
\textbf{Date} & \textbf{Test} & \textbf{QI} & \textbf{N} & \textbf{Contenu KB} \\
\midrule
20 Avr & H7 Issues seules & 54 & 50 & Issues H7 uniquement \\
20 Avr & H7 Prompt Fix & 57 & 50 & + System prompt optimisé \\
24 Avr & H7 Config Tool Test & 61 & 50 & + Hybrid search calibré \\
28 Avr & H7 + Diagnostic Cards & 94 & 50 & + Resolver Cases + Diag \\
08 Mai & H7 Real Cases ST Community & 88 & 4 & Test sur cas réels forum \\
02 Juin & Test élargi & 83 & 4 & Multi-séries first attempt \\
08 Juin & Multi-séries (H7+F4+H5+U5+WL) & 54 & 52 & 5 séries complètes \\
\bottomrule
\end{longtable}

La Figure~\ref{fig:qi_evolution} visualise cette évolution, mettant en évidence le saut majeur apporté par les Diagnostic Cards et la régression lors du passage multi-séries.

% Graphique QI evolution - exporter depuis draw.io ou Excel
\begin{figure}[H]
\centering
\includegraphics[width=0.9\textwidth]{qi_evolution_chart.png}
\caption{Évolution du Quality Index au fil des améliorations incrémentales (draw.io)}
\label{fig:qi_evolution}
\end{figure}

\textit{Le graphique en barres montre l'évolution : Issues seules (54) $\to$ +Prompt (57) $\to$ +Hybrid (61) $\to$ +Diag Cards (94, $\uparrow$+33) $\to$ Real Cases (88) $\to$ Multi 4Q (83) $\to$ Multi 52Q (54, dilution multi-séries). Le saut majeur de 61 à 94 démontre l'impact des Diagnostic Cards.}

\subsection{Analyse de l'Évolution}

\subsubsection{Phase 1 : Issues seules (QI = 54)}

Le score initial de 54 avec les issues H7 seules montre que les issues structurées (symptôme, root cause, fix) fournissent déjà une base solide pour le RAG. Les lacunes :
\begin{itemize}
  \item Pas de documentation de fichiers → questions sur les API non couvertes
  \item Pas de chaînes Issue→PR→Commit → pas de traçabilité du fix
  \item Pas de cartes de diagnostic → les problèmes complexes multi-étapes mal couverts
\end{itemize}

\subsubsection{Phase 2 : Ajout des Resolver Cases et Diagnostic Cards (QI = 94)}

Le saut de 61 à \textbf{94} est le gain le plus important du projet. Les \textbf{Diagnostic Cards} structurent l'information sous forme :

\begin{quote}
Symptôme → Hypothèses → Validation → Root Cause → Fix → Prévention
\end{quote}

Cette structure correspond exactement à la façon dont un ingénieur raisonne face à un problème embedded. Le retrieval RAG trouve ces cartes car elles contiennent des \textbf{query aliases} (reformulations multiples du problème) et des \textbf{keywords techniques} (noms d'API, registres, erreurs).

Les \textbf{Resolver Cases} ajoutent la dimension de traçabilité :
\begin{quote}
Issue \#42 → PR \#87 → Commit abc123 → Fichier modifié : \texttt{stm32f4xx\_hal\_spi.c}
\end{quote}

\subsubsection{Phase 3 : Passage multi-séries (QI = 54)}

Le score chute à 54 lors de l'ajout de 4 nouvelles séries (F4, H5, U5, WL). Cette chute s'explique par :

\begin{enumerate}
  \item \textbf{Dilution du contexte} : la KB passe de $\sim$2\,000 à $\sim$14\,000 documents. Les queries retournent des documents de séries différentes, diluant la pertinence.
  \item \textbf{Questions nouvelles} : le jeu de test passe de 50 à 52 questions, avec 10 questions sur les nouvelles séries pour lesquelles les Diagnostic Cards ne sont pas encore aussi mûres.
  \item \textbf{Conflit de noms} : des API identiques entre séries (ex: \texttt{HAL\_SPI\_Init} existe dans F4 et H7) créent de l'ambiguïté lors du retrieval.
\end{enumerate}

\textbf{Leçon} : le passage à l'échelle nécessite un \textbf{filtrage par série} plus agressif dans les métadonnées des documents, ou un prompt qui pousse la persona à inclure le nom de la série dans ses queries KB.

\subsection{Raisonnement de Chaque Amélioration}

\begin{longtable}{p{4cm}p{5cm}p{4cm}}
\toprule
\textbf{Amélioration} & \textbf{Raisonnement} & \textbf{Impact QI} \\
\midrule
System prompt optimisé & Forcer la recherche KB avant toute réponse & +3 (54→57) \\
Hybrid search & Combiner semantic + full-text pour ne rater ni les reformulations ni les identifiants & +4 (57→61) \\
Diagnostic Cards & Structurer les problèmes complexes en cartes raisonnées & +33 (61→94) \\
Resolver Cases & Ajouter la traçabilité Issue→PR→Fix & Inclus dans +33 \\
Alfred restrictif & Empêcher le court-circuit de la KB & Récupéré $-10$ points \\
Files V3 (nettoyage) & Réduire le bruit des CMSIS headers & +3 estimé \\
PDF descriptions & Rendre les figures PDF cherchables & +1 estimé \\
\bottomrule
\end{longtable}

\section{Couverture Actuelle et Future}
\label{sec:couverture}

\subsection{13 Séries en Production}

Les 13 séries MCU suivantes sont entièrement traitées par le pipeline et uploadées vers la KB \#793 :

\begin{longtable}{lcccccc}
\toprule
\textbf{Série} & \textbf{Repos} & \textbf{Issues} & \textbf{Files} & \textbf{Resolver} & \textbf{Diag.} & \textbf{Alfred} \\
\midrule
STM32CubeC0 & 4 & \checkmark & \checkmark & --- & \checkmark & \checkmark \\
STM32CubeF0 & 5 & \checkmark & \checkmark & --- & \checkmark & \checkmark \\
STM32CubeF1 & 6 & \checkmark & \checkmark & \checkmark & \checkmark & \checkmark \\
STM32CubeF2 & 4 & \checkmark & \checkmark & --- & \checkmark & \checkmark \\
STM32CubeF3 & 5 & \checkmark & \checkmark & --- & \checkmark & \checkmark \\
STM32CubeF4 & 16 & \checkmark & \checkmark & \checkmark & \checkmark & \checkmark \\
STM32CubeF7 & 8 & \checkmark & \checkmark & \checkmark & \checkmark & \checkmark \\
STM32CubeG0 & 5 & \checkmark & \checkmark & --- & \checkmark & \checkmark \\
STM32CubeH5 & 6 & \checkmark & \checkmark & \checkmark & \checkmark & \checkmark \\
STM32CubeH7 & 12 & \checkmark & \checkmark & \checkmark & \checkmark & \checkmark \\
STM32CubeH7RS & 4 & \checkmark & \checkmark & --- & \checkmark & \checkmark \\
STM32CubeU5 & 8 & \checkmark & \checkmark & \checkmark & \checkmark & \checkmark \\
STM32CubeWL & 5 & \checkmark & \checkmark & \checkmark & \checkmark & \checkmark \\
\bottomrule
\end{longtable}

\textit{Note : ``---'' pour Resolver signifie que les issues de cette série n'ont pas de chaînes Issue$\to$PR$\to$Commit complètes, pas que le pipeline ne supporte pas le format.}

\subsection{4 Nouvelles Séries Exportées (Pipeline Terminé)}

Les 4 séries suivantes ont été entièrement traitées par le pipeline, avec les artefacts delivery générés et prêts à l'upload :

\begin{longtable}{lccccc}
\toprule
\textbf{Série} & \textbf{Issues} & \textbf{Files} & \textbf{Resolver} & \textbf{Diag. Cards} & \textbf{Statut} \\
\midrule
STM32CubeG4 & 89 & 3\,875 & 0 & 4 & Prêt à uploader \\
STM32CubeL0 & 23 & 2\,614 & 2 & 3 & Prêt à uploader \\
STM32CubeL1 & 11 & 2\,847 & 0 & 3 & Prêt à uploader \\
STM32CubeL4 & 68 & 5\,231 & 1 & 8 & Prêt à uploader \\
\bottomrule
\end{longtable}

\textbf{Observations} :
\begin{itemize}
  \item \textbf{G4 (0 Resolver)} : les 89 issues G4 ont été fermées sans PRs formelles liées. Les corrections ont été intégrées dans des releases globales sans lien explicite issue$\to$PR.
  \item \textbf{L4 (5\,231 files)} : c'est la série avec le plus de fichiers parmi les nouvelles, reflétant un écosystème BSP riche (NUCLEO, Discovery, Eval boards).
  \item \textbf{Alfred Images} : 25 images échouées (API 404 intermittent) sur G4, L0, L4. Corrigible par re-run du script d'enrichissement.
\end{itemize}

\subsection{8 Séries Configurées (Prêtes à Lancer)}

Les configurations (fichiers \texttt{config\_*.json}) sont créées pour : L5, N6, U0, U3, WB, WB0, WBA, WL3. Le déploiement nécessite :
\begin{enumerate}
  \item Exécution du pipeline : \texttt{python -m pipeline.run\_full\_workflow --export-st-ready}
  \item Création des datasources sur la KB \#793
  \item Upload via les scripts d'automatisation
\end{enumerate}
